% TOP_PROBLEM: {"id": 30007159, "title": "Spectral scalar-curvature control of codimension-two width", "category_id": 6, "category": {"id": 6, "name": "geometry", "display_name": "Geometry", "description": "Euclidean and non-Euclidean geometry, geometric structures.", "slug": "geometry", "order_index": 6, "created_at": "2026-07-31T15:26:25.671Z"}, "status": "open"}
% FIRST_POSTED: 2026-10-06
% ATTEMPT_STATUS: solved
% OPENAI_CLAIM: {"schema_version":1,"source_repo":"https://github.com/openai/math","source_commit":"adc7f1241b42e322a6451854ab7e4b4c146bf78a","release_date":"2026-10-06","match":"full","resolution":"proved","independently_reviewed":false,"second_pass":"agreed","families":[{"family":"336","title":"Spectral scalar curvature, Urysohn width, and macroscopic dimension","subject":"Differential geometry","manuscripts":[{"title":"Spectral scalar curvature and uniform Urysohn width","pdf_path":"preprints/Spectral-scalar-curvature-and-uniform-Urysohn-width-October-5-2026/main.pdf","readme_path":"preprints/Spectral-scalar-curvature-and-uniform-Urysohn-width-October-5-2026/README.md","date":"2026-10-05","theorem_ref":"thm:main (sections/introduction.tex:35)","statement_sources":[{"path":"preprints/Spectral-scalar-curvature-and-uniform-Urysohn-width-October-5-2026/build/sections/introduction.tex","line":35,"label":"thm:main"}]}],"lean":null,"reasoning_trace":null}]}
\documentclass[11pt]{article}
\usepackage[margin=25mm]{geometry}
\usepackage{fontspec}
\setmainfont{Latin Modern Roman}
\usepackage{amsmath,amssymb,mathtools}
\usepackage{unicode-math}
\setmathfont{Latin Modern Math}
\usepackage{enumitem}
\usepackage{xurl,hyperref,bookmark}
\hypersetup{colorlinks=true,urlcolor=blue}
\setcounter{secnumdepth}{0}
\setlength{\parindent}{0pt}
\setlength{\parskip}{5pt}
\setlength{\emergencystretch}{3em}
\begin{document}
\section*{Spectral scalar-curvature control of codimension-two width}

\subsection{OpenAI's claim}
\textit{Higher-dimensional spectral width}: The spectral lower bound $-4\Delta+\mathrm{Scal}\ge1$ is claimed to imply uniform continuous $(n-2)$-width for every complete $n$-manifold with $n\ge4$. The bound uses the initial metric rather than an auxiliary conformal distance.

\subsection{Mathematical statement of the claimed result}
\paragraph{Mathematical claim 1.} For $n\ge4$, a connected complete smooth boundaryless $(M^n,g)$ satisfying \[\int_M(4|\nabla\phi|^2+\operatorname{Scal}_g\phi^2)\,dV_g\ge\int_M\phi^2\,dV_g\quad(\phi\in C_c^\infty(M;\mathbb R))\] is claimed to have a continuous $f:M\to K$ with $\dim K\le n-2$ and all fiber diameters at most $C_n$, depending only on $n$. Thus its Urysohn $(n-2)$-width, the infimum of these uniform fiber-diameter bounds, is at most $C_n$.

\subsection{Outline of the argument}
\textit{Higher-dimensional spectral width}: A positive supersolution produces a signed drift with a weighted curvature bound. An auxiliary graph metric and updated drift add a quadratic reserve even when the original drift is unbounded. The subsequent stretching and wall construction uses that reserve to preserve curvature while reducing intersection multiplicity, yielding fibers bounded in the original geometry.

\subsection{Scope relative to this problem}
\textit{Higher-dimensional spectral width}: The spectral lower bound $-4\Delta+\mathrm{Scal}\ge1$ is claimed to imply uniform continuous $(n-2)$-width for every complete $n$-manifold with $n\ge4$. The bound uses the initial metric rather than an auxiliary conformal distance. The parameters, regularity assumptions and normalization in this question are retained in applying that assertion.

With the hypotheses and conventions above, the claimed result addresses the whole mathematical assertion on this page.

\subsection{Formal verification}
\textit{Higher-dimensional spectral width}: The pinned catalogue lists no Lean formalization for this result.

This site has not run the Lean development. These links describe OpenAI’s selected formal statements and should be read with their stated scope.

\subsection{Status caveat}
OpenAI reports that an internal model produced the manuscripts. Their arguments have not been independently verified as proofs by this site. Any portion outside the documented formal scope may contain mistakes, and the catalogue label reports the announced mathematical claim.

\subsection{Sources}
\begin{itemize}
\item OpenAI, \emph{Spectral scalar curvature and uniform Urysohn width}, 2026-10-05; citation key \texttt{OAI:Spectral-scalar-curvature-and-uniform-Urysohn-width-October-5-2026}. \href{https://github.com/openai/math/blob/adc7f1241b42e322a6451854ab7e4b4c146bf78a/preprints/Spectral-scalar-curvature-and-uniform-Urysohn-width-October-5-2026/main.pdf}{PDF (GitHub)}; \href{https://github.com/openai/math/raw/adc7f1241b42e322a6451854ab7e4b4c146bf78a/preprints/Spectral-scalar-curvature-and-uniform-Urysohn-width-October-5-2026/main.pdf}{PDF (direct)}; \href{https://github.com/openai/math/blob/adc7f1241b42e322a6451854ab7e4b4c146bf78a/preprints/Spectral-scalar-curvature-and-uniform-Urysohn-width-October-5-2026/README.md}{Manuscript page and citation}; \href{https://github.com/openai/math/blob/main/preprints/Spectral-scalar-curvature-and-uniform-Urysohn-width-October-5-2026/main.pdf}{Latest version}; Statement source: \href{https://github.com/openai/math/blob/adc7f1241b42e322a6451854ab7e4b4c146bf78a/preprints/Spectral-scalar-curvature-and-uniform-Urysohn-width-October-5-2026/build/sections/introduction.tex\#L35}{thm:main}.
\item \href{https://github.com/openai/math/blob/adc7f1241b42e322a6451854ab7e4b4c146bf78a/CONTENTS.md}{OpenAI catalogue entry}.
\end{itemize}
\end{document}
